\section{完备性让逼近留在空间内}
\label{sec:12-Real-Analysis-Support/03-Topology-Basics/03}

在开区间 \(X=(0,1)\) 上定义一个迭代：
\[
T(x)=\frac{x}{2}.
\]
从 \(x_0=0.8\) 开始：\(0.8, 0.4, 0.2, 0.1, 0.05,\dots\) 每一步减半，无限逼近 0。数字越来越小，彼此间距越来越紧密——它们明显``应该''收敛。但 0 不在 \((0,1)\) 里。这个序列表现出了收敛的一切外在条件，极限却空降到空间外面去了。

这不是一个病态构造。多项式序列
\[
p_n(x)=\sum_{k=0}^n\frac{x^k}{k!}
\]
在 \([0,1]\) 上一致收敛到 \(e^x\)。每个 \(p_n\) 都是多项式，而极限 \(e^x\) 不是多项式。Weierstrass 逼近定理甚至宣告：\(C([a,b])\) 中\emph{每一个}连续函数都是多项式序列的一致极限。多项式空间一次次逼近成功，极限却住在隔壁的连续函数空间里。

问题出在哪？不是极限不存在，是空间本身漏掉了一些``应该属于它的点''。

\begin{definition}
度量空间 \((X,d)\) 称为\textbf{完备的}，如果 \(X\) 中每个 Cauchy 数列都收敛到 \(X\) 中的某个点。

数列 \((x_n)\) 是 Cauchy 的，若对任意 \(\varepsilon>0\)，存在 \(N\) 使 \(m,n\ge N\) 时 \(d(x_m,x_n)<\varepsilon\)。
\end{definition}

Cauchy 条件捕捉的是``项与项之间无限靠近''，比``存在极限''更可操作：你不需要预先猜出极限是谁。完备性保证的正是：只要项与项之间无限靠近，极限就\emph{一定}在空间里。

\subsection{闭子集继承完备性，但反过来要小心}

若 \(X\) 完备、\(F\subseteq X\) 闭，那么 \(F\) 也完备。任何 \(F\) 中的 Cauchy 列在 \(X\) 中收敛（因 \(X\) 完备）；极限点落在 \(F\) 的闭包里，而 \(F\) 闭，所以极限就在 \(F\) 里。

反过来，完备空间中的完备子集一定是闭的。若 \(F\) 不闭，必存在 \(F\) 中序列收敛到 \(X\setminus F\) 的点；这个序列是 Cauchy 的，极限却不在 \(F\) 中，与 \(F\) 的完备性矛盾。

这解释了为什么许多迭代定理要求在\emph{闭}可行集上工作。即使环境空间完备，开子集仍可能让极限落到边界外——开头 \(T(x)=x/2\) 在 \((0,1)\) 上的失败，正是开集完备性缺失的直接后果。

注意完备性和闭性的区别：闭性问的是``子集在环境空间中是否包含自己的极限点''；完备性问的是``空间自身是否包含所有 Cauchy 逼近的终点''。前者依赖环境空间，后者是内在性质。

\subsection{函数空间的完备性}

在 \(C([a,b])\) 上赋予 sup 范数
\[
\norm{f}_\infty = \sup_{x\in[a,b]}|f(x)|,
\]
这个空间是完备的。若 \((f_n)\) 在 sup 范数下是 Cauchy 列，逐点极限 \(f(x)=\lim_n f_n(x)\) 存在（因为对每个 \(x\)，\((f_n(x))\) 是实数的 Cauchy 列，实数完备）。从一致 Cauchy 性出发：对任意 \(\varepsilon>0\)，存在 \(N\) 使所有 \(m,n\ge N\) 时 \(\sup_x|f_m(x)-f_n(x)|<\varepsilon\)。固定 \(m\ge N\)，令 \(n\to\infty\)，得 \(\sup_x|f_m(x)-f(x)|\le\varepsilon\)，即 \(f_n\to f\) 一致。连续函数的一致极限仍连续（上一单元已证），所以 \(f\in C([a,b])\)。

完备赋范线性空间称为 Banach 空间。\(C([a,b])\)、\(L^p\) 空间以及大量数值方法的解空间都在这个框架内。

相比之下，多项式空间在 sup 范数下不完备：Cauchy 列可以一致收敛到非多项式函数，极限始终属于更大的 \(C([a,b])\)——空间太小，装不下自己的逼近产物。

\subsection{紧致与完备：两个相似但不同的概念}

紧致度量空间一定完备。Cauchy 列有收敛子列（序列紧致），而 Cauchy 性迫使全序列趋向同一极限：\(d(x_n,x)\le d(x_n,x_{n_k})+d(x_{n_k},x)\)，第一项由 Cauchy 条件控制，第二项由子列收敛控制。

完备\emph{不}推出紧致。\(\R\) 完备但无界（你可以一直往前走）。甚至在闭且有界的范围内，无限维空间也可能不紧。度量空间中，紧致的精确刻画是：
\[
\text{完备} \;+\; \text{全有界} \;=\; \text{紧致}.
\]
全有界意味着对任意 \(\varepsilon>0\)，集合能被有限个半径为 \(\varepsilon\) 的球覆盖。

普通有界只说集合落在某个大球里；全有界要求在\emph{任意精细}的尺度上都能用有限网格兜住。以 \(\ell^2\) 为例（平方可和序列空间），标准基向量 \(e_i\) 全部落在单位球内，且任意两个不同的基向量距离为 \(\sqrt{2}\)。取半径小于 \(\sqrt{2}/2\) 的球，每个球至多装入一个 \(e_i\)；有限个球永远盖不住无穷多个相互远离的基向量。所以无限维闭单位球不全有界，因而也不紧——这正是上一节 Heine--Borel 失效的根本原因。

\subsection{Banach 不动点定理：完备性发挥计算威力}

完备性的计算价值在 Banach 不动点定理中展现得最充分。

设 \(T:X\to X\) 是\emph{压缩映射}：
\[
d(Tx,Ty)\le q\,d(x,y),\qquad 0<q<1.
\]
若 \(X\) 非空且完备，则：

\begin{theorem}[Banach 不动点定理]
在上述条件下：
\begin{enumerate}
\tightlist
\item \(T\) 有唯一不动点 \(x^\star\)；
\item 任取 \(x_0\in X\)，迭代 \(x_{n+1}=T(x_n)\) 都收敛到 \(x^\star\)；
\item 误差有几何上界：
  \[
  d(x_n,x^\star)\le\frac{q^n}{1-q}\,d(x_1,x_0).
  \]
\end{enumerate}
\end{theorem}

推导过程即是一座小型的``完备性应用示范''。首先由压缩性：
\[
d(x_{n+1},x_n)\le q^n\,d(x_1,x_0).
\]
对 \(m>n\) 累加并用几何级数求和：
\[
d(x_m,x_n)\le\bigl(q^{m-1}+\cdots+q^n\bigr)\,d(x_1,x_0)
\le\frac{q^n}{1-q}\,d(x_1,x_0).
\]
右端随 \(n\to\infty\) 趋零，\((x_n)\) 是 Cauchy 列。完备性保证极限 \(x^\star\) 在 \(X\) 内。令 \(m\to\infty\) 即得几何误差界。压缩条件同时说明 \(T\) 连续（它是 Lipschitz 常数 \(q<1\) 的 Lipschitz 映射），因此在 \(x_{n+1}=T(x_n)\) 两侧取极限得 \(T(x^\star)=x^\star\)。

唯一性：若有两个不动点 \(x^\star,y^\star\)，则
\[
d(x^\star,y^\star)=d(Tx^\star,Ty^\star)\le q\,d(x^\star,y^\star),
\]
迫使 \(d(x^\star,y^\star)=0\)。

回到 \((0,1)\) 上的 \(T(x)=x/2\)：它是压缩映射，\(q=1/2\)，迭代逼近 0，但 \((0,1)\) 不完备（它缺了 0），因此 Banach 定理的前提不成立——空间验证这一步绝对不能略过。

\subsection{完备性的边界：不缺任何 Cauchy 列的终点}

Banach 定理支撑不动点迭代、常微分方程局部存在唯一性、积分方程和动态规划。每次应用都要求三项核对：空间完备；映射把空间映入自身；压缩常数严格小于 1 且与起始点无关。

完备性回答的是一个基础问题：你的逼近过程会不会把答案推导到空间外面去？如果答案可能在空间外，那么收敛性证明的第一步就通不过——极限根本无处安放。

从压缩迭代到函数空间的极限，从数值方法到概率论中的 \(L^2\) 完备化，完备性始终扮演同一个角色：你的 Cauchy 逼近的终点，已经有位置在等你。
